\documentclass[11pt,letterpaper]{article}
\usepackage[margin=1in]{geometry}
\usepackage{amsmath,amssymb,amsthm}
\usepackage{physics}
\usepackage{hyperref}
\usepackage{cleveref}
\usepackage{booktabs}

\newtheorem{theorem}{Theorem}[section]
\newtheorem{lemma}[theorem]{Lemma}
\newtheorem{proposition}[theorem]{Proposition}
\newtheorem{corollary}[theorem]{Corollary}
\newtheorem{definition}[theorem]{Definition}

\numberwithin{equation}{section}
\renewcommand{\theequation}{S6.\arabic{equation}}

\title{Supplement S6: Neutrino Mass Derivation \\
\large Supporting Manuscript \S13}
\author{UDT Collaboration}
\date{}

\begin{document}
\maketitle

\begin{abstract}
We derive the neutrino mass from the UDT metric via the full chain:
action integral, electromagnetic vertex, $\alpha^3$ factorization, and the
mass formula $m_\nu = \alpha^3 m_e/(2j+1)^2$. The $\alpha^3$ factor
arises from a triple application of the electromagnetic vertex (the neutrino
as a grade-6 object constructed from three grade-2 forms). The $1/(2j+1)^2 = 1/4$
coefficient comes from the spin-squared normalization. We prove the orbit
matching theorem on the 7-chain and derive cross-check ratios. Verified by
\texttt{analysis/10\_wedge\_product.py}.
\end{abstract}

\tableofcontents

%----------------------------------------------------------------------
\section{The neutrino in the exterior algebra}
\label{sec:exterior}
%----------------------------------------------------------------------

Recall from Supplement~S4 that the Diophantine identity selects a
9-dimensional base space $\mathbb{R}^9$ with $\dim = 2l + 2|\kappa_\mathrm{max}| + 1 = 9$.
The exterior algebra $\Lambda^*(\mathbb{R}^9)$ organizes particles by grade:

\begin{center}
\begin{tabular}{cccl}
\toprule
Grade $k$ & $\dim\Lambda^k$ & Particle & Role \\
\midrule
0 & 1 & vacuum & identity \\
1 & 9 & --- & base vectors \\
2 & 36 & $\alpha$-boson & electromagnetic coupling \\
3 & 84 & pion & strong sector \\
4 & 126 & --- & (transition grade) \\
5 & 126 & --- & (transition grade) \\
6 & 84 & neutrino & Hodge dual of pion \\
7 & 36 & --- & dual of $\alpha$ \\
8 & 9 & --- & dual of base \\
9 & 1 & --- & top form \\
\bottomrule
\end{tabular}
\end{center}

The neutrino lives at grade~6, which is the Hodge dual of grade~3
(the pion). The Hodge star on $\Lambda^9$:
\begin{equation}
\label{eq:hodge}
  *: \Lambda^3(\mathbb{R}^9) \xrightarrow{\sim} \Lambda^6(\mathbb{R}^9),
  \qquad \dim\Lambda^3 = \dim\Lambda^6 = 84.
\end{equation}

\subsection{Triple wedge construction}

The grade-6 form can be constructed as a triple wedge product of grade-2 forms:
\begin{equation}
\label{eq:triple_wedge}
  \omega_\nu = \alpha \wedge \alpha \wedge \alpha \in \Lambda^6(\mathbb{R}^9),
\end{equation}
where $\alpha \in \Lambda^2(\mathbb{R}^9)$ is the electromagnetic 2-form
(36 components). This construction gives the neutrino its $\alpha^3$
coupling: each wedge factor contributes one power of the fine structure
constant.

%----------------------------------------------------------------------
\section{The $\alpha^3$ factorization}
\label{sec:alpha3}
%----------------------------------------------------------------------

\subsection{Action integral}

The neutrino effective action on the UDT background involves the
triple-vertex diagram where the fermion line (electron) emits and
reabsorbs three virtual photons, generating the neutrino mass:
\begin{equation}
\label{eq:action}
  \mathcal{A}_\nu = \int\!\left(\bar{\psi}\gamma^\mu A_\mu\right)^3\,
  e^{3\phi}\,r^2\,dr\,d\Omega.
\end{equation}

Each vertex contributes a factor $\sqrt{\alpha}$ (coupling constant), and
the three vertices contribute $\alpha^{3/2}$. Squaring the amplitude to
get the mass shift:
\begin{equation}
  m_\nu \propto |\mathcal{A}_\nu|^2 / m_e \propto \alpha^3 \cdot m_e.
\end{equation}

\subsection{Detailed factorization}

The fine structure constant decomposes as (Supplement~S5):
\begin{equation}
  \alpha = \frac{I_2}{36\pi} = \frac{I_2}{\binom{9}{2}\,\pi}.
\end{equation}

Therefore:
\begin{equation}
\label{eq:alpha3}
  \alpha^3 = \left(\frac{I_2}{36\pi}\right)^3 = \frac{I_2^3}{36^3\,\pi^3}
  = \frac{I_2^3}{46656\,\pi^3}.
\end{equation}

Note that $36^3 = 46656 = \dim\left(\Lambda^2\right)^{\otimes 3}$, the full
(unconstrained) tensor product dimension. The antisymmetrization to
$\Lambda^6$ reduces this by a factor $46656/84 = 555.43\ldots$

%----------------------------------------------------------------------
\section{The mass formula}
\label{sec:mass_formula}
%----------------------------------------------------------------------

\begin{theorem}[Neutrino mass]
\label{thm:neutrino}
\begin{equation}
\label{eq:m_nu}
\boxed{
  m_\nu = \frac{\alpha^3\,m_e}{(2j+1)^2} = \frac{\alpha^3\,m_e}{4}.
}
\end{equation}
\end{theorem}

\subsection{The $1/(2j+1)^2$ coefficient}

The denominator $(2j+1)^2 = 4$ arises from the spin-squared normalization
of the neutrino vertex. Unlike the pion (which couples through
$(2j+1)^2(2l+1)(2|\kappa_\mathrm{max}|+1) = 84$ channels), the neutrino
couples through the \emph{inverse} of the spin-squared factor because the
Hodge dual reverses the role of spin and orbit:

\begin{itemize}
  \item Pion (grade~3): coefficient includes $(2j+1)^2$ in the \emph{numerator}.
  \item Neutrino (grade~6 = Hodge dual of grade~3): coefficient has
    $(2j+1)^2$ in the \emph{denominator}.
\end{itemize}

This Hodge duality between grades 3 and 6 is responsible for the neutrino's
extreme lightness: it is suppressed by both $\alpha^3$ and $1/4$ relative
to the electron.

\subsection{Numerical evaluation}

Using $\alpha_\mathrm{PDG} = 1/137.036$ and $m_e = 0.51100$~MeV:
\begin{align}
  m_\nu &= \frac{(1/137.036)^3 \times 0.51100}{4}~\mathrm{MeV} \nonumber\\
  &= \frac{3.886 \times 10^{-7} \times 0.51100}{4}~\mathrm{MeV} \nonumber\\
  &= 4.97 \times 10^{-8}~\mathrm{MeV} \nonumber\\
  &= 0.0497~\mathrm{eV}. \label{eq:m_nu_num}
\end{align}

This should be compared with $\sqrt{\Delta m^2_{31}} = \sqrt{2.515 \times 10^{-3}}~\mathrm{eV} = 0.0501$~eV from neutrino oscillation data (NuFIT~5.2, normal ordering), yielding an agreement of $\sim 1\%$.

%----------------------------------------------------------------------
\section{Orbit matching theorem}
\label{sec:orbit_matching}
%----------------------------------------------------------------------

\subsection{The 7-chain}

The orbit space $\mathbb{R}^{2|\kappa_\mathrm{max}|+1} = \mathbb{R}^7$ can
be endowed with a graph structure: the \emph{7-chain} is the path graph
$P_7$ with 7 vertices and 6 edges:
\begin{equation}
  v_1 - v_2 - v_3 - v_4 - v_5 - v_6 - v_7.
\end{equation}

\begin{definition}[Near-perfect matching]
A \emph{near-perfect matching} on a graph with an odd number of vertices is
a matching that covers all but one vertex. On $P_7$, a near-perfect matching
consists of 3 disjoint edges.
\end{definition}

\begin{theorem}[Orbit matching]
\label{thm:orbit}
The 7-chain $P_7$ has exactly 4 near-perfect matchings.
\end{theorem}

\textit{Proof.}
A near-perfect matching on $P_7$ leaves exactly one unmatched vertex. The
unmatched vertex can be:
\begin{itemize}
  \item $v_1$: matching $\{(v_2,v_3), (v_4,v_5), (v_6,v_7)\}$,
  \item $v_3$: matching $\{(v_1,v_2), (v_4,v_5), (v_6,v_7)\}$,
  \item $v_5$: matching $\{(v_1,v_2), (v_3,v_4), (v_6,v_7)\}$,
  \item $v_7$: matching $\{(v_1,v_2), (v_3,v_4), (v_5,v_6)\}$.
\end{itemize}
These are the only possibilities: the unmatched vertex must have even index
from each end to allow consecutive-edge matchings on a path graph. On $P_7$,
the valid unmatched vertices are $v_1, v_3, v_5, v_7$ (the odd-indexed
vertices). \qed

\subsection{Connection to $1/(2j+1)^2$}

The 4 near-perfect matchings correspond to the $(2j+1)^2 = 4$ spin channels.
Each matching represents one way the neutrino's internal structure can be
arranged on the 7-dimensional orbit, and the $1/4$ in the mass formula
counts the average over these matchings.

%----------------------------------------------------------------------
\section{Cross-check ratios}
\label{sec:cross_checks}
%----------------------------------------------------------------------

\subsection{$m_p \cdot m_\nu / m_e^2$}

\begin{equation}
  \frac{m_p \cdot m_\nu}{m_e^2}
  = \frac{(6\pi^5 m_e)(\alpha^3 m_e/4)}{m_e^2}
  = \frac{3\pi^5\alpha^3}{2}.
\end{equation}

Numerically: $\frac{3 \times 305.97 \times 3.886 \times 10^{-7}}{2} = 1.78 \times 10^{-4}$.

This ratio involves only the angular quantum numbers (through $\pi^5$ and $\alpha$),
providing a cross-check that connects the heaviest baryon to the lightest lepton.

\subsection{$m_\mu / m_\nu$}

\begin{equation}
  \frac{m_\mu}{m_\nu}
  = \frac{20\pi^3 m_e/3}{\alpha^3 m_e/4}
  = \frac{80\pi^3}{3\alpha^3}.
\end{equation}

Numerically: $\frac{80 \times 31.006}{3 \times 3.886 \times 10^{-7}} = 2.13 \times 10^9$.

The enormous ratio $m_\mu/m_\nu \sim 10^9$ is generated by $1/\alpha^3 \sim 10^{6.4}$
combined with the $\pi^3$ factor, without any hierarchy problem---both masses are
derived from the same quantum numbers.

%----------------------------------------------------------------------
\section{Damping ratio}
\label{sec:damping}
%----------------------------------------------------------------------

The ratio of consecutive exterior algebra dimensions provides the
\emph{damping ratio}:
\begin{equation}
\label{eq:damping}
  \frac{\dim\Lambda^3}{\dim\Lambda^4} = \frac{\binom{9}{3}}{\binom{9}{4}}
  = \frac{84}{126} = \frac{2}{3}.
\end{equation}

And the even/odd ratio:
\begin{equation}
\label{eq:even_odd}
  \frac{\dim\Lambda^2}{\dim\Lambda^3} = \frac{\binom{9}{2}}{\binom{9}{3}}
  = \frac{36}{84} = \frac{3}{7}.
\end{equation}

These ratios govern the transition amplitudes between grades in the exterior
algebra and constrain the neutrino coupling structure.

%----------------------------------------------------------------------
\section{Why $\alpha^3$ and not $\alpha^2$ or $\alpha^4$}
\label{sec:why_alpha3}
%----------------------------------------------------------------------

The power of $\alpha$ is fixed by the grade structure:
\begin{itemize}
  \item The neutrino lives at grade~6 in $\Lambda^*(\mathbb{R}^9)$.
  \item The electromagnetic 2-form lives at grade~2.
  \item The number of 2-forms needed to reach grade~6 is $6/2 = 3$.
  \item Therefore the neutrino requires $\alpha^{6/2} = \alpha^3$.
\end{itemize}

No other power is consistent with the grade assignment. An $\alpha^2$
neutrino would live at grade~4 (which has $\dim = 126$, not matching the
pion's Hodge dual), and an $\alpha^4$ neutrino would require grade~8
($\dim = 9$, too small).

%----------------------------------------------------------------------
\section{Verification scripts}
\label{sec:verification}
%----------------------------------------------------------------------

\begin{itemize}
  \item \texttt{analysis/10\_wedge\_product.py}: Exterior algebra structure,
    Hodge duality, damping ratios, Euler characteristic.
  \item Neutrino mass computed in \texttt{lib/angular\_integrals.py}
    function \texttt{neutrino\_mass()}.
  \item Output: \texttt{data/generated/10\_wedge\_product.json}.
\end{itemize}

\end{document}
